%! The Tangent Function — Unit 7, Lesson 7.1 — Student Packet Cover Sheet
\documentclass[10pt]{article}
\usepackage{precalculus-article}
\usepackage{precalculus-boxes}
\usepackage{ltablex}
\keepXColumns

\begin{document}

% Full-bleed navy banner
\begin{tikzpicture}[remember picture, overlay]
  \fill[navy] (current page.north west)
    rectangle ([yshift=-0.9in]current page.north east);
\end{tikzpicture}
\vspace{-0.8in}

\begin{center}
  {\color{white}\textbf{\LARGE Precalculus}}\\[4pt]
  {\color{skymid}\large\textbf{Unit 7: Analytic Trigonometry and Polar Coordinates}}\\[6pt]
  {\color{white}\textbf{\large Lesson 7.1 \quad The Tangent Function}}
\end{center}

\vspace{0.22in}
\namedateperiod
\vspace{0.18in}

\begin{learningtargetbox}
\small
\begin{enumerate}[leftmargin=1.5em, itemsep=5pt]
  \item I can construct a representation of $f(\theta)=\tan\theta$ using the unit
    circle by computing $\sin\theta/\cos\theta$ and interpreting it as the slope of
    the terminal ray. \hfill\textit{(LO 3.8.A)}
  \item I can describe key characteristics of the tangent function: period $\pi$,
    location of vertical asymptotes at $\theta=\pi/2+k\pi$, the function is always
    increasing, and the graph changes from concave down to concave up within each
    period. \hfill\textit{(LO 3.8.B)}
  \item I can describe the effect of additive and multiplicative transformations on
    the tangent function, including changes to period, phase shift, vertical dilation,
    and vertical translation. \hfill\textit{(LO 3.8.C)}
\end{enumerate}
\end{learningtargetbox}

\vspace{0.14in}

\begin{tocbox}
\small
\renewcommand{\arraystretch}{1.5}
\begin{tabularx}{\linewidth}{@{} c l X r @{}}
\rowcolor{sky}
\textbf{\#} & \textbf{Component} & \textbf{Description} & \textbf{Score} \\
\midrule
1 & Warm-Up        & Spiral review: unit circle values, slope, and asymptotes & \blank{1.2cm} \\
2 & Guided Notes   & Tangent from the unit circle; period, asymptotes, transformations & \blank{1.2cm} \\
3 & Group Activity & Tiered tangent practice (R/A/E) & \blank{1.2cm} \\
4 & Exit Ticket    & Compute $\tan\theta$; identify period and asymptote & \blank{1.2cm} \\
5 & Homework       & Mixed practice + AP-style MC + preview of Lesson 3.9 & \blank{1.2cm} \\
\midrule
  & \multicolumn{2}{r}{\textbf{Total}} & \blank{1.2cm} \\
\end{tabularx}
\end{tocbox}

\end{document}
